\SourcePage{400}
\SGAsection{3}{射影丛的 $K'$ 的结构定理的推论}

\begin{SGAresult}{命题}{3.1}
设 $S$ 为\Term{quasi-compact}{拟紧}\Term{scheme}{概形}，$\mathcal E_1,\ldots,\mathcal E_n$ 为\Term{finite-type}{有限型}\Term{locally-free-sheaf}{局部自由} $\mathcal O_S$-\Term{module}{模}，$\mathrm D(\mathcal E_i)$ 为 $\mathcal E_i$ 的\Term{flag-scheme}{旗概形}，并令
\[
X=\mathrm D(\mathcal E_1)\times_S\cdots\times_S\mathrm D(\mathcal E_n).
\]
则\Term{canonical}{典范}\Term{homomorphism}{同态} $K'(S)\longrightarrow K'(X)$ 是\Term{injectivity}{单射}。这里 $K'$ 既可由\Term{finite-type}{有限型}\Term{locally-free-sheaf}{局部自由} $\mathcal O_S$-\Term{module}{模}定义，也可由\Term{perfect-complex}{完美复形}定义。
\end{SGAresult}
\begin{proof}
利用\Term{flag-scheme}{旗概形}的\Term{universal-property}{泛性质}，立即可知 $X$ 可以如下构造：考虑一列 $S$-\Term{scheme}{概形} $X_1,\ldots,X_n$，递归地取 $X_k$ 为 $\mathcal E_k$ 在 $X_{k-1}$ 上的\Term{inverse-image}{逆像}的\Term{flag-scheme}{旗概形}；于是 $X=X_n$。因此，为证明命题，只须证明 $n=1$ 的情形，\Term{generic}{一般}情形即可由归纳法得到。故可设 $X$ 为某个\Term{finite-type}{有限型}\Term{locally-free-sheaf}{局部自由} $\mathcal O_S$-\Term{module}{模} $\mathcal E$ 的\Term{flag-scheme}{旗概形}。

已知这时 $X$ 可以如下构造：置 $\mathcal E=\mathcal F_0$，$S=X_0$，递归地定义\Term{scheme}{概形} $X_1,\ldots,X_n$ 以及对每个 $i$ 的\Term{finite-type}{有限型}\Term{locally-free-sheaf}{局部自由} $\mathcal O_{X_i}$-\Term{module}{模} $\mathcal F_i$，令
\SourcePage{401}
\[
X_k=\mathbf P(\mathcal F_{k-1}),\qquad
\mathcal F_k=\operatorname{Ker}\bigl(f_k^*(\mathcal F_{k-1})\longrightarrow\mathcal O_{X_k}(1)\bigr),
\]
其中 $f_k$ 为\Term{structure-morphism}{结构态射} $X_k\longrightarrow X_{k-1}$。

于是只须证明 $X$ 为 $S$ 上的\Term{projective-bundle}{射影丛}时的命题，\Term{flag-bundle}{旗丛}的情形即可由归纳法\Term{pushout}{推出}。而这一情形直接来自 1.1（相应地，2.1）。
\end{proof}

\begin{SGAresult}{\Term{theorem}{定理}}{3.2}
设 $S$ 为\Term{quasi-compact}{拟紧}\Term{scheme}{概形}。由\Term{finite-type}{有限型}\Term{locally-free-sheaf}{局部自由} $\mathcal O_S$-\Term{module}{模}定义的\Term{pre-lambda-ring}{预 $\lambda$-环} $K'(S)$ 是一个\Term{lambda-ring}{$\lambda$-环}（第 V 讲 2.4）。
\end{SGAresult}
\begin{proof}
由于\Term{homomorphism}{同态} $\lambda_t\colon K'(S)\longrightarrow 1+K'(S)[[t]]^+$ 是\Term{additive}{加性}的，为验证它是一个\Term{lambda-homomorphism}{$\lambda$-同态}，只须对 $\operatorname{cl}(\mathcal E)$ 型的元素验证，其中 $\mathcal E$ 为\Term{finite-type}{有限型}\Term{locally-free-sheaf}{局部自由} $\mathcal O_S$-\Term{module}{模}。因而只须证明：对 $K'(S)$ 中任意两个分别为\Term{finite-type}{有限型}\Term{locally-free-sheaf}{局部自由} $\mathcal O_S$-\Term{module}{模} $\mathcal E,\mathcal F$ 的类的元素 $E,F$，第 V 讲 2.4.2 的关系成立：
\[
\begin{aligned}
\lambda^i(EF)&=P_i\bigl(\lambda^1(E),\ldots,\lambda^i(E);\lambda^1(F),\ldots,\lambda^i(F)\bigr),\\
\lambda^j(\lambda^i(E))&=Q_{i,j}\bigl(\lambda^1(E),\ldots,\lambda^{ij}(E)\bigr).
\end{aligned}
\]
由于 $S$ \Term{quasi-compact}{拟紧}，可用\Term{finite}{有限}个两两不交的\Term{open-subset}{开集}\Term{covering}{覆盖} $S$，使 $\mathcal E$ 与 $\mathcal F$ 在其中每个\Term{open-subset}{开集}上均为常\Term{rank}{秩}。而\Term{finite}{有限}个\Term{open-subset}{开集}的\Term{disjoint-union}{不交并}的 $K'$ 是这些\Term{open-subset}{开集}的 $K'$ 的\Term{direct-product}{直积}。因此只须在这个\Term{covering}{覆盖}中各个\Term{open-subset}{开集}的 $K'$ 中验证第 V 讲 2.4.2 的关系。故可假定 $\mathcal E$ 与 $\mathcal F$ 在 $S$ 上均为常\Term{rank}{秩}。
\SourcePage{402}
令 $X=\mathrm D(\mathcal E)\times_S\mathrm D(\mathcal F)$。由 3.1，\Term{lambda-homomorphism}{$\lambda$-同态} $K'(S)\longrightarrow K'(X)$ 是\Term{injectivity}{单射}。因此，为证明第 V 讲 2.4.2 的关系，可转到 $K'(X)$ 中。在 $K'(X)$ 中，元素 $E,F$ \Term{resolution}{分解}为\Term{invertible-sheaf}{可逆层}的类之和，故由\Term{additive}{加性}，只须对\Term{invertible-sheaf}{可逆层}的类 $\xi,\eta$ 证明这些关系；将其写成第 V 讲 2.4.1 的\Term{form}{形式}，即为
\[
1+\xi\eta t=(1+\xi t)\mathbin{\widetilde\otimes}(1+\eta t),\qquad
\lambda_u(1+\xi t)=1+\{1+\xi t\}u.
\]
这正是定义 $1+K'(S)[[t]]^+$ 的\Term{lambda-ring}{$\lambda$-环}结构的第 V 讲 2.3.1、2.3.3 两个关系。
\end{proof}

现指出如下更\Term{generic}{一般}的结果。
\begin{SGAresult}{\Term{theorem}{定理}}{3.3}
设 $(T,A)$ 为一个交换\Term{ringed-topos}{环化拓扑斯}；以 $K$ 表示 $T$ 上\Term{rank}{秩}\Term{bounded}{有界}的\Term{locally-free-sheaf}{局部自由} $A$-\Term{module}{模}的\Term{grothendieck-group}{Grothendieck 群}。则 $K$ 配备其\Term{canonical}{典范}\Term{pre-lambda-structure}{预 $\lambda$-结构}（第 V 讲 2.2 b)）后，是一个\Term{lambda-ring}{$\lambda$-环}。
\end{SGAresult}
\begin{proof}
为证明 $K$ 是\Term{lambda-ring}{$\lambda$-环}，须证明第 V 讲 2.4.1 的关系；由\Term{additive}{加性}，只须对分别为\Term{rank}{秩}\Term{bounded}{有界}的\Term{locally-free-sheaf}{局部自由} $A$-\Term{module}{模}之类的 $E,F$ 证明。由于 $\mathcal E,\mathcal F$ 的\Term{rank}{秩}\Term{bounded}{有界}，可将 $T$ \Term{resolution}{分解}为\Term{finite}{有限}个\Term{open-subset}{开集}的“\Term{disjoint-union}{不交并}”，使 $\mathcal E,\mathcal F$ 在各个\Term{open-subset}{开集}上均为常\Term{rank}{秩}。利用 $K'$ 与\Term{finite}{有限}\Term{disjoint-union}{不交并}的相容性，并将第 V 讲 2.4.1 的关系写成第 V 讲 2.4.2 的\Term{form}{形式}，可知只须在上述各\Term{open-subset}{开集}的 $K'$ 中证明第 V 讲 2.4.2 的关系。换言之，只须对分别为常\Term{rank}{秩}\Term{locally-free-sheaf}{局部自由} $A$-\Term{module}{模} $\mathcal E,\mathcal F$ 的类的元素 $E,F$ 证明这些关系。
\SourcePage{403}
设 $\mathcal E$ 为\Term{rank}{秩} $n$ 的\Term{locally-free-sheaf}{局部自由} $A$-\Term{module}{模}。存在一个 $\mathrm{GL}(n,A)$ 下的\Term{torsor}{挠子} $P$，使
\[
\mathcal E=P\mathbin{\times^{\mathrm{GL}(n,A)}}A^n,
\]
其中通过标准\Term{representation}{表示}将 $A^n$ 看作 $\mathrm{GL}(n,A)$-\Term{module}{模}。于是可定义从 $R_{\mathbf Z}(\mathrm{GL}(n))$ 到 $K$ 的\Term{lambda-homomorphism}{$\lambda$-同态}：将 $\mathrm{GL}(n)$ 在\Term{rank}{秩} $p$ 的自由 $\mathbf Z$-\Term{module}{模}中的任一\Term{representation}{表示} $\rho$，对应到\Term{locally-free-sheaf}{局部自由} $A$-\Term{module}{模} $P\mathbin{\times^{\mathrm{GL}(n,A)}}A^p$ 的类，这里通过 $\rho$ 将 $A^p$ 看作 $\mathrm{GL}(n,A)$-\Term{module}{模}。在这一\Term{homomorphism}{同态}下，$\mathcal E$ 的类就是标准\Term{representation}{表示}的像。因此，为验证 $\lambda^i(E)$ 之间的一个\Term{polynomial}{多项式}关系，只须验证 $R_{\mathbf Z}(\mathrm{GL}(n))$ 中的标准\Term{representation}{表示}满足这一关系。

对\Term{rank}{秩}分别为 $n,p$ 的两个\Term{locally-free-sheaf}{局部自由} $A$-\Term{module}{模} $\mathcal E,\mathcal F$，可作类似论证。此时可找到 $\mathrm{GL}(n,A)\times\mathrm{GL}(p,A)$ 下的\Term{torsor}{挠子} $P$，使
\[
\mathcal E=P\mathbin{\times^{\mathrm{GL}(n,A)\times\mathrm{GL}(p,A)}}A^n,
\]
其中 $\mathrm{GL}(n,A)$ 通过标准\Term{representation}{表示}\Term{action}{作用}，$\mathrm{GL}(p,A)$ 通过平凡\Term{representation}{表示}\Term{action}{作用}，由此将 $A^n$ 看作 $\mathrm{GL}(n,A)\times\mathrm{GL}(p,A)$-\Term{module}{模}；$\mathcal F$ 的情形反之。再以与上面相同的方式定义从 $R_{\mathbf Z}(\mathrm{GL}(n)\times\mathrm{GL}(p))$ 到 $K$ 的\Term{lambda-homomorphism}{$\lambda$-同态}，于是归结为对 $R_{\mathbf Z}(\mathrm{GL}(n)\times\mathrm{GL}(p))$ 中相应的元素验证 $\lambda^i(E),\lambda^j(F)$ 之间的关系。

故只须证明：对 $\operatorname{Spec}\mathbf Z$ 上任一由\Term{finite}{有限}个\Term{generic}{一般}\Term{linear}{线性}\Term{group-scheme}{群概形}之积构成的\Term{group-scheme}{群概形} $G$，\Term{pre-lambda-ring}{预 $\lambda$-环} $R_{\mathbf Z}(G)$ 是\Term{lambda-ring}{$\lambda$-环}。为此可引用 [5]，其中显式计算了 $R_{\mathbf Z}(G)$。如不使用这一结果，也可利用 1.14，按与 3.2 类似的方式论证。
\end{proof}

\SourcePage{404}
\SGAsection{4}{旗丛的 $K'$ 的计算}
\textbf{4.1。} 设 $S$ 为\Term{scheme}{概形}，$\mathcal E$ 为\Term{rank}{秩} $n$ 的\Term{locally-free-sheaf}{局部自由} $\mathcal O_S$-\Term{module}{模}。设 $\pi=(p_1,\ldots,p_k)$ 为一列不小于 $1$ 的整数，满足
\[
\sum_{i=1}^k p_i=n.
\]
回顾一下，$\mathcal E$ 的 $\pi$ 型\Term{flag-functor}{旗函子}是指 $S$-\Term{scheme}{概形}\Term{category}{范畴}上的如下\Term{set}{集合}值\Term{functor}{函子}：它将一个 $S$-\Term{scheme}{概形} $f\colon X\longrightarrow S$ 对应到 $f^*(\mathcal E)$ 的\Term{quotient}{商}列
\[
(\mathcal L_{k-1},\ldots,\mathcal L_1)
\]
所构成的\Term{set}{集合}，其中对每个 $i$，$\mathcal L_i$ 是 $\mathcal L_{i+1}$ 的\Term{quotient}{商}，且为\Term{rank}{秩} $p_1+\cdots+p_i$ 的\Term{locally-free-sheaf}{局部自由层}。

这个\Term{functor}{函子}可由一个 $S$-\Term{scheme}{概形} $\mathrm D_\pi(\mathcal E)$ \Term{representable-functor}{表示}，它带有 $\mathcal E$ 的\Term{inverse-image}{逆像}的一列\Term{canonical}{典范}\Term{quotient}{商}，它们\Term{locally-free-sheaf}{局部自由}，\Term{rank}{秩}依次为 $p_1,p_1+p_2,\ldots$。

若 $\pi=(1,n-1)$，则 $\mathrm D_\pi(\mathcal E)=\mathbf P(\mathcal E)$。

若 $\pi=(r,n-r)$，则 $\mathrm D_\pi(\mathcal E)=\mathrm G_{r,n}(\mathcal E)$，即 $\mathcal E$ 的 $(r,n-r)$ 型\Term{grassmannian}{Grassmann 概形}。

若所有 $p_i$ 均等于 $1$，则得到 $\mathcal E$ 的通常\Term{flag-bundle}{旗丛}，记作 $\mathrm D(\mathcal E)$。

\begin{SGAresult}{\Term{lemma}{引理}}{4.2}
设 $S$ 为\Term{scheme}{概形}，$\mathcal E$ 为\Term{rank}{秩} $n$ 的\Term{locally-free-sheaf}{局部自由} $\mathcal O_S$-\Term{module}{模}，$\pi=(p_1,\ldots,p_i,\ldots,p_k)$，且 $p_1+\cdots+p_k=n$。设 $p',p''$ 为两个不小于 $1$ 的整数，满足 $p'+p''=p_i$；置
\[
\pi'=(p_1,\ldots,p_{i-1},p',p'',p_{i+1},\ldots,p_k),\qquad \pi''=(p',p'').
\]
设 $\mathcal L_1,\ldots,\mathcal L_{k-1}$ 为 $\mathcal E$ 在 $\mathrm D_\pi(\mathcal E)$ 上的\Term{inverse-image}{逆像}的\Term{canonical}{典范}\Term{quotient}{商}，$\mathcal F$ 为 $\mathcal L_i\longrightarrow\mathcal L_{i-1}$ 的\Term{kernel}{核}。则有\Term{canonical}{典范}的 $S$-\Term{isomorphism}{同构}
\SourcePage{405}
\[
\mathrm D_{\pi'}(\mathcal E)\simeq\mathrm D_{\pi''}(\mathcal F).
\]
\end{SGAresult}
\begin{proof}
由\Term{flag-scheme}{旗概形}的\Term{functorial-definition}{函子定义}立即可得，留给读者。
\end{proof}

\begin{SGAresult}{\Term{lemma}{引理}}{4.3}
设 $A$ 为有\Term{identity-element}{单位元}的\Term{commutative-ring}{交换环}，$n>0$ 为整数，$(c_i)_{1\leq i\leq n}$ 为 $A$ 中的一族元素；设 $C$ 为由\Term{generator}{生成元} $\xi_i$（$i=1,\ldots,n$）在关系
\[
\sigma_i(\xi_1,\ldots,\xi_n)=c_i,\qquad i=1,\ldots,n
\]
下生成的 $A$-\Term{algebra}{代数}，其中 $\sigma_i$ 表示初等\Term{symmetric-function}{对称函数}。设 $p_1,\ldots,p_k$ 为不小于 $1$、且和为 $n$ 的整数，并对 $1\leq j\leq k$ 令
\[
\xi^{(j)}=(\xi_1^{(j)},\ldots,\xi_{p_j}^{(j)}),\qquad
\xi_i^{(j)}=\xi_{\sum_{\alpha<j}p_\alpha+i}.
\]
置
\[
c_i^{(j)}=\sigma_i(\xi_1^{(j)},\ldots,\xi_{p_j}^{(j)}),\qquad
1\leq j\leq k,\quad 1\leq i\leq p_j,
\]
并设 $B$ 为由这些 $c_i^{(j)}$ 生成的 $C$ 的\Term{subalgebra}{子代数}。则：

\textup{i)} 形如
\[
(\xi^{(1)})^{\alpha(1)}\cdots(\xi^{(k)})^{\alpha(k)},
\qquad
\alpha(j)=(\alpha_1^{(j)},\ldots,\alpha_{p_j}^{(j)}),\quad
0\leq\alpha_i^{(j)}\leq p_j-i
\]
的元素构成 $C$ 作为 $B$-\Term{module}{模}的一组基。

\textup{ii)} $A$-\Term{algebra}{代数} $B$ 的\Term{generator}{生成元} $c_i^{(j)}$ 满足从关系
\[
\sum_{i=1}^n c_i=\prod_{j=1}^k\left(\sum_{i=1}^{p_j}c_i^{(j)}\right)
\]
\Term{pushout}{推出}的 $n$ 个等权关系（赋予 $c_i,c_i^{(j)}$ 的权为 $i$），而且这 $n$ 个关系生成 $c_i^{(j)}$ 之间的关系\Term{ideal}{理想}。
\end{SGAresult}
\SourcePage{406}
\begin{SGAresult}{推论}{4.4}
$C$ 中形如
\[
\xi_1^{\alpha_1}\cdots\xi_{n-1}^{\alpha_{n-1}},\qquad 0\leq\alpha_i\leq n-i
\]
的元素构成 $C$ 作为 $A$-\Term{module}{模}的一组基。
\end{SGAresult}

这个\Term{lemma}{引理}及其推论取自 [2]，第 4 讲，第 4、19 页；证明亦见该处。

\textbf{4.5。} 设 $S$ 为\Term{quasi-compact}{拟紧}\Term{scheme}{概形}，$\mathcal E$ 为\Term{rank}{秩} $n$ 的\Term{locally-free-sheaf}{局部自由} $\mathcal O_S$-\Term{module}{模}，$\pi=(p_1,\ldots,p_k)$ 为一列不小于 $1$、且和为 $n$ 的整数。置 $X=\mathrm D_\pi(\mathcal E)$；$X$ 带有一列\Term{canonical}{典范}\Term{quotient}{商} $\mathcal L_i$；置 $\mathcal L_0=0$，$\mathcal L_k=f^*(\mathcal E)$，其中 $f$ 为\Term{structure-morphism}{结构态射} $X\longrightarrow S$。对 $1\leq i\leq k$，以 $\mathcal F_i$ 表示 $\mathcal L_i\longrightarrow\mathcal L_{i-1}$ 的\Term{kernel}{核}。设 $K'(X)$ 为 $X$ 上\Term{finite-type}{有限型}\Term{locally-free-sheaf}{局部自由层}（相应地，\Term{perfect-complex}{完美复形}）的\Term{grothendieck-group}{Grothendieck 群}。以 $F_j$ 表示 $\mathcal F_j$ 的类，以 $\lambda_i^{(j)}$ 表示 $\Lambda^i(\mathcal F_j)$ 的类。在用\Term{perfect-complex}{完美复形}定义 $K'(X)$ 的情形，级数 $\lambda_t(E),\lambda_t(F_j)$ 仍有显然的意义，其中 $E$ 为 $\mathcal E$ 的类。最后注意
\[
\sum_{j=1}^k F_j=E.
\]

\begin{SGAresult}{\Term{theorem}{定理}}{4.6}
在 4.5 的假设与记号下，$K'(S)$-\Term{algebra}{代数} $K'(X)$ 由 $\lambda_i^{(j)}$ 生成；它们之间的关系\Term{ideal}{理想}由下式\Term{pushout}{推出}的关系生成：
\[
\lambda_t(E)=\prod_{j=1}^k\lambda_t(F_j).
\]
\SourcePage{407}
换言之，$K'(X)$ 是由\Term{generator}{生成元} $F_1,\ldots,F_k$ 在下列关系下生成的 $K'(S)$-\Term{lambda-algebra}{$\lambda$-代数}（参见第 V 讲 4.6）：
\[
\textup{i)}\quad \lambda^p(F_i)=0\quad(p>p_i),\qquad
\textup{ii)}\quad \sum_i F_i=E.
\]
\end{SGAresult}

先考虑所有 $p_i$ 均等于 $1$ 的情形，此时\Term{theorem}{定理}表述如下。
\begin{SGAresult}{推论}{4.7}
设 $S$ 为\Term{quasi-compact}{拟紧}\Term{scheme}{概形}，$\mathcal E$ 为\Term{rank}{秩} $n$ 的\Term{locally-free-sheaf}{局部自由} $\mathcal O_S$-\Term{module}{模}，$X$ 为 $\mathcal E$ 的\Term{flag-bundle}{旗丛}，$\xi_1,\ldots,\xi_n$ 为\Term{invertible-sheaf}{可逆层} $\mathcal F_1,\ldots,\mathcal F_n$ 的类（3.5 的记号）。则 $K'(X)$ 是由 $\xi_i$ 在关系
\[
\sigma_i(\xi_1,\ldots,\xi_n)=\lambda^i(E)
\]
下生成的 $K'(S)$-\Term{algebra}{代数}。
\end{SGAresult}
\begin{proof}
对 $n$ 作归纳。$n=1$ 时\Term{theorem}{定理}显然成立。设已对 $n-1$ 证明了\Term{theorem}{定理}；考虑 $\mathbf P(\mathcal E)$ 及其上的\Term{exact-sequence}{正合列}
\[
0\longrightarrow\mathcal F\longrightarrow g^*(\mathcal E)
\longrightarrow\mathcal O_{\mathbf P(\mathcal E)}(1)\longrightarrow0,
\]
其中 $g$ 为 $\mathbf P(\mathcal E)$ 的\Term{structure-morphism}{结构态射}。由\Term{lemma}{引理} 4.2，$X$ 等同于 $\mathbf P(\mathcal E)$ 上的通常\Term{flag-bundle}{旗丛} $\mathrm D(\mathcal F)$，且这一等同对应各个\Term{canonical}{典范}\Term{sheaf}{层}。图
\[
\begin{tikzcd}
\mathrm D(\mathcal E)=\mathrm D(\mathcal F)\arrow[r,"h"]\arrow[dr,"f"'] & \mathbf P(\mathcal E)\arrow[d,"g"]\\
&S
\end{tikzcd}
\]
表明 $K'(S)\longrightarrow K'(X)$ 经 $K'(\mathbf P(\mathcal E))$ \Term{resolution}{分解}。将归纳假设用于 $\mathrm D(\mathcal F)$，并由推论 3.4，$K'(X)$ 在 $K'(\mathbf P(\mathcal E))$ 上有一组基，其元素形如
\[
\xi_2^{\alpha_2}\cdots\xi_{n-1}^{\alpha_{n-1}},\qquad 0\leq\alpha_i\leq n-i;
\]
另一方面，$K'(\mathbf P(\mathcal E))$ 在 $K'(S)$ 上以 $1,\xi_1,\ldots,\xi_1^{n-1}$ 为基。因此 $K'(X)$ 在 $K'(S)$ 上有一组基，其元素为
\SourcePage{408}
\[
\xi_1^{\alpha_1}\cdots\xi_{n-1}^{\alpha_{n-1}},\qquad 0\leq\alpha_i\leq n-i.
\]
又由 $\lambda_t$ 的\Term{additive}{加性}，$\prod_i(1+\xi_it)=\lambda_t(E)$，故推论中的关系成立。

设 $C$ 是由\Term{generator}{生成元} $\xi_i$ 在关系
\[
\sigma_i(\xi_1,\ldots,\xi_n)=\lambda^i(E)
\]
下生成的 $K'(S)$-\Term{algebra}{代数}，则存在一个从 $C$ 到 $K'(X)$ 的满\Term{homomorphism}{同态}。由 4.4，这个\Term{homomorphism}{同态}将 $C$ 的一组基映成 $K'(X)$ 的一组基，因而是\Term{isomorphism}{同构}，推论得证。
\end{proof}

\begin{proof}[\Term{theorem}{定理} 4.6 的证明]
现处理 $\pi$ 型\Term{flag-bundle}{旗丛}的\Term{generic}{一般}情形。通常\Term{flag-bundle}{旗丛} $\mathrm D(\mathcal E)$ 可从 $\mathrm D_\pi(\mathcal E)$ 出发，依次构造各 $\mathcal F_j$ 的\Term{flag-bundle}{旗丛}而得（4.2）。由此，$K'(\mathrm D(\mathcal E))$ 作为 $K'(X)$-\Term{module}{模}，有一组由 4.3 i) 中的元素组成的基；这里 $\xi_i^{(j)}$ 是在 $\mathrm D(\mathcal E)$ 上\Term{splitting}{分裂}各 $\mathcal F_j$ 的\Term{inverse-image}{逆像}时所出现的、类似于 $\mathcal F_j$ 的那些\Term{sheaf}{层}的类（3.5 和 3.7 的记号）。故 $K'(X)$ 是 $K'(\mathrm D(\mathcal E))$ 的\Term{subring}{子环}。此外，$\lambda_t$ 的\Term{additive}{加性}给出
\[
\lambda_i^{(j)}=\sigma_i(\xi_1^{(j)},\ldots,\xi_{p_j}^{(j)}),
\]
而 $K'(X)$ 包含 $K'(\mathrm D(\mathcal E))$ 中由 $\lambda_i^{(j)}$ 生成的 $K'(S)$-\Term{subalgebra}{子代数}。

以 $B$ 表示这个\Term{subalgebra}{子代数}。由 4.3 i) 及刚才的论述，$K'(\mathrm D(\mathcal E))$ 中有一族元素，它既是其 $B$-\Term{module}{模}结构的一组基，又是其 $K'(X)$-\Term{module}{模}结构的一组基；从而 $K'(X)=B$，这证明了 $K'(X)$ 由 $\lambda_i^{(j)}$ 生成。由 4.3 ii)，再结合 4.7 所确定的 $K'(\mathrm D(\mathcal E))$ 在 $K'(S)$ 上的结构，以及 $\xi_i^{(j)}$ 是\Term{splitting}{分裂} $\mathcal E$ 在 $\mathrm D(\mathcal E)$ 上的\Term{inverse-image}{逆像}时所出现的\Term{sheaf}{层}的类这一事实，即得 $\lambda_i^{(j)}$ 之间的关系\Term{ideal}{理想}。
\SourcePage{409}
这就证明了 4.6 的第一种\Term{form}{形式}。

为\Term{pushout}{推出}第二种\Term{form}{形式}，注意 $K'(S)$-\Term{lambda-algebra}{$\lambda$-代数} $K'(X)$ 中的元素 $F_i$ 满足关系 i)、ii)。设 $A$ 是由满足 i)、ii) 的元素 $F'_i$ 生成的 $K'(S)$-\Term{lambda-algebra}{$\lambda$-代数}，则有一个 $K'(S)$-\Term{lambda-homomorphism}{$\lambda$-同态} $A\longrightarrow K'(X)$，将 $F'_i$ 映到 $F_i$（第 V 讲 4.7）。由 4.6，这个\Term{homomorphism}{同态}是\Term{surjectivity}{满射}。由于 $K'(S)[F_1,\ldots,F_k]$ 中由关系 i)、ii) 生成的\Term{lambda-ideal}{$\lambda$-理想}包含由 4.6 中各关系生成的通常\Term{ideal}{理想}，这个\Term{homomorphism}{同态}也是\Term{injectivity}{单射}，故为\Term{isomorphism}{同构}。
\end{proof}

\begin{SGAresult}{推论}{4.8}
在 4.5 的假设与记号下，置
\[
c_i^{(j)}=\gamma^i(F_j-p_j),\qquad 1\leq i\leq p_j.
\]
则 $K'(S)$-\Term{algebra}{代数} $K'(X)$ 由 $c_i^{(j)}$ 生成；它们之间的关系\Term{ideal}{理想}由下式\Term{pushout}{推出}的关系生成：
\[
\gamma_t(E-n)=\prod_{j=1}^k\gamma_t(F_j-p_j).
\]
\end{SGAresult}
\begin{proof}
这只是 4.6 的另一种表述。事实上，$c_i^{(j)}$ 可以用 $\lambda_i^{(j)}$ 表示，反之亦然。而关系 $\lambda_t(E)=\prod_j\lambda_t(F_j)$ \Term{equivalence}{等价}于关系 $\gamma_t(E-n)=\prod_j\gamma_t(F_j-p_j)$。
\end{proof}

\textbf{4.9。} 设 $S$ 为\Term{scheme}{概形}；以 $\mathrm G_S^{n,r}$ 表示 $(\mathcal O_S)^r$ 的 $(n,r-n)$ 型\Term{grassmannian}{Grassmann 概形}，其中 $r\geq n$。当 $r$ 变化时，各 $\mathrm G_S^{n,r}$ 构成一个\Term{inductive-system}{归纳系}；\Term{morphism}{态射} $\mathrm G_S^{n,r}\longrightarrow\mathrm G_S^{n,r+1}$ 的定义是将 $(\mathcal O_{\mathrm G_S^{n,r}})^r$ 看作 $(\mathcal O_{\mathrm G_S^{n,r}})^{r+1}$ 对第 $(r+1)$ 个因子取的\Term{quotient}{商}。
\SourcePage{410}
因此，在这个\Term{morphism}{态射}系中，各\Term{canonical}{典范}\Term{quotient}{商}在\Term{inverse-image}{逆像}下相互对应。将如此定义的\Term{inductive-system}{归纳系}记为 $\mathrm G_S^{n,\infty}$。

由公式
\[
K'(\mathrm G_S^{n,\infty})=\varprojlim_r K'(\mathrm G_S^{n,r})
\tag{4.9.1}
\]
定义它的 $K'$。

设 $\mathcal F_1,\mathcal F_2$ 为 $\mathrm G_S^{n,r}$ 上的\Term{canonical}{典范}\Term{sheaf}{层}（4.5），$F_1,F_2$ 为它们在 $K'(\mathrm G_S^{n,r})$ 中的类。置 $c_i=\gamma^i(F_1-n)$。当 $r$ 变化时，各 $c_i$ 在过渡\Term{morphism}{态射}下相互对应，因而定义了 $K'(\mathrm G_S^{n,\infty})$ 中的元素，仍记为 $c_i$。

\begin{SGAresult}{命题}{4.10}
设 $S$ 为\Term{quasi-compact}{拟紧}\Term{scheme}{概形}。则在 4.9 的记号下，有\Term{isomorphism}{同构}
\[
K'(\mathrm G_S^{n,\infty})\simeq K'(S)[[c_1,\ldots,c_n]].
\]
\end{SGAresult}
\begin{proof}
由 4.8，$K'(\mathrm G_S^{n,r})$ 是由 $c_i=\gamma^i(F_1-n)$（$1\leq i\leq n$）与 $\gamma^j(F_2-r+n)$（$1\leq j\leq r-n$）在从下式\Term{pushout}{推出}的关系下生成的 $K'(S)$-\Term{algebra}{代数}：
\[
\gamma_t(F_1-n)\gamma_t(F_2-r+n)=\gamma_t(r-r)=1,
\]
亦即
\[
\gamma_t(F_2-r+n)=1/\gamma_t(F_1-n).
\]
这个关系可将 $\gamma^j(F_2-r+n)$ 用 $c_i$ 表示，并通过令 $1/\gamma_t(F_1-n)$ 在\Term{degree}{次数} $>r-n$ 的系数为零，给出 $c_i$ 之间一系列等权关系（赋予 $c_i$ 的权为 $i$）；而且这些关系生成 $c_i$ 之间的关系\Term{ideal}{理想}。因此可写成
\SourcePage{411}
\[
K'(\mathrm G_S^{n,r})=K'(S)[c_1,\ldots,c_n]/I_r,
\]
其中 $I_r$ 是 $K'(S)[c_1,\ldots,c_n]$ 中一个\Term{degree}{次数} $>r-n$ 的\Term{homogeneous-ideal}{齐次理想}。于是由定义 (4.9.1)，$K'(\mathrm G_S^{n,\infty})$ 是 $K'(S)[c_1,\ldots,c_n]$ 关于各 $I_r$ 所定义的\Term{topology}{拓扑}的\Term{separated-completion}{分离完备化}。这个\Term{topology}{拓扑}比按权定义的\Term{topology}{拓扑}细；我们将证明两者\Term{equivalence}{等价}，由此即得命题。

置 $A=K'(S)[c_1,\ldots,c_n]/I_r$，并设 $C$ 为由\Term{generator}{生成元} $x_1,\ldots,x_n$ 在关系
\[
\sigma_j(x_1,\ldots,x_n)=c_j
\]
下生成的 $A$-\Term{algebra}{代数}，其中 $\sigma_j$ 表示初等\Term{symmetric-function}{对称函数}。由 4.4，\Term{canonical}{典范}\Term{homomorphism}{同态} $A\longrightarrow C$ 是\Term{injectivity}{单射}。可得
\[
\gamma_t(F_1-n)=1+c_1t+\cdots+c_nt^n=\prod_{i=1}^n(1+x_it).
\]
由于 $1/\gamma_t(F_1-n)$ 是\Term{polynomial}{多项式}，每个 $1/(1+x_it)$ 也是\Term{polynomial}{多项式}。因此各 $x_i$ 均\Term{nilpotent}{幂零}，且存在整数 $N$，使任何以 $x_i$ 为变量、以 $K'(S)$ 中元素为系数、\Term{degree}{次数} $>N$ 的\Term{polynomial}{多项式}均为零。因而 $A$ 中权 $>N$ 的\Term{polynomial}{多项式}均为零，即 $K'(S)[c_1,\ldots,c_n]$ 中权 $>N$ 的\Term{polynomial}{多项式}均属于 $I_r$。因此所考虑的两个\Term{topology}{拓扑}\Term{equivalence}{等价}，证明完毕。
\end{proof}

\SGAsection{5}{在射影丛上的应用；$f_*$ 的研究}
\textbf{5.1。} 本节中，$S$ 为\Term{quasi-compact}{拟紧}\Term{scheme}{概形}，$\mathcal E$ 为\Term{rank}{秩} $r+1$ 的\Term{locally-free-sheaf}{局部自由} $\mathcal O_S$-\Term{module}{模}，$X$ 为与 $\mathcal E$ 相伴的\Term{projective-bundle}{射影丛}。在 $X$ 上有\Term{canonical}{典范}\Term{exact-sequence}{正合列}
\[
0\longrightarrow\mathcal F\longrightarrow f^*(\mathcal E)\longrightarrow\mathcal O_X(1)\longrightarrow0,
\]
\SourcePage{412}
其中 $f$ 为\Term{structure-morphism}{结构态射} $X\longrightarrow S$。

这里所考虑的 $K'$ 均由\Term{finite-type}{有限型}\Term{locally-free-sheaf}{局部自由层}定义。以 $E,F,\xi$ 分别表示 $\mathcal E,\mathcal F,\mathcal O_X(1)$ 的类。置 $x=1-\xi^{-1}$，$x'=\xi-1=\xi x$，并以 $\bar x,\bar x'$ 表示 $x,x'$ 在 $\operatorname{Gr}^1K'(X)$ 中的类。由于 $\xi-1$ 与 $x$ 均属于 $\operatorname{Fil}^1K'(X)$，有 $(\xi-1)x\in\operatorname{Fil}^2K'(X)$，即
\[
x\equiv\xi x=x'\pmod{\operatorname{Fil}^2K'(X)},\qquad\text{亦即}\quad\bar x'=-\bar x.
\]

\textbf{5.2。} 若 $\mathcal G$ 为\Term{finite-type}{有限型}\Term{locally-free-sheaf}{局部自由} $\mathcal O_X$-\Term{module}{模}，则各 $R^if_*(\mathcal G)$ \Term{generic}{一般}并不是 $S$ 上的\Term{finite-type}{有限型}\Term{locally-free-sheaf}{局部自由层}，因而不能像用\Term{perfect-complex}{完美复形}定义 $K'$ 时那样定义一个\Term{homomorphism}{同态} $f_*\colon K'(X)\longrightarrow K'(S)$。借助\Term{theorem}{定理} 1.1，可以如下避开这个困难：若 $y\in K'(X)$，则可以唯一地写成
\[
y=a_0+a_1\xi^{-1}+\cdots+a_r\xi^{-r},\qquad a_i\in K'(S);
\tag{5.2.1}
\]
于是置
\[
f_*(y)=a_0.
\tag{5.2.2}
\]
如此定义的\Term{homomorphism}{同态}显然具有\Term{additive}{加性}，并满足\Term{projection-formula}{投影公式}，即为 $K'(S)$-\Term{linear}{线性}的。

可验证以下关系：
\[
f_*(1)=1;\qquad f_*(\xi^n)=0\quad(-r\leq n<0);\qquad
f_*(\xi^n)=\sigma_n\quad(1\leq n).
\tag{5.2.3}
\]
为验证最后一组关系，对 $n$ 作归纳。设对 $p\leq n-1$，有 $f_*(\xi^p)=\sigma_p$；将关系 (1.9.1) 乘以 $\xi$ 的适当幂，得
\[
\sum_{i=0}^{r+1}(-1)^i\lambda^i(E)\xi^{n-i}=0.
\]
再施以 $f_*$，得
\SourcePage{413}
\[
\sum_{i=0}^{r+1}(-1)^i\lambda^i(E)f_*(\xi^{n-i})=0.
\]
由关系 (1.11.3) 及归纳假设立即得到所需结果。

\begin{SGAresult}{命题}{5.3}
在 4.1 的记号与假设下，对每个 $k\geq0$，有如下 $K'(S)$-\Term{module}{模}的\Term{direct-sum-decomposition}{直和分解}：
\[
\operatorname{Fil}^kK'(X)=\operatorname{Fil}^kK'(S)+\operatorname{Fil}^{k-1}K'(S)\,x+\cdots+\operatorname{Fil}^{k-r}K'(S)\,x^r,
\]
这里所考虑的\Term{filtration}{滤过}是第 V 讲 3.9.1 和 3.10 中定义的\Term{augmented-algebra}{增广} \Term{lambda-algebra}{$\lambda$-代数}的\Term{filtration}{滤过}。
\end{SGAresult}
\begin{proof}
首先注意 $x$ 的\Term{augmentation}{增广}为零。此外有
\[
\gamma_t(x)=\lambda_{t/(1-t)}(1-\xi^{-1})
=\left(1+\frac t{1-t}\right)\left(1+\frac{\xi^{-1}t}{1-t}\right)^{-1}
=\frac1{1-xt}.
\]
因此，对每个 $i$，$\gamma^i(x)=x^i$。从而 $K'(X)$ 是由\Term{generator}{生成元} $X$ 在对一切 $i$ 的关系 $\gamma^i(X)=X^i$ 下生成的 $K'(S)$-\Term{lambda-algebra}{$\lambda$-代数}的一个\Term{quotient}{商}，取\Term{quotient}{商}所用的\Term{ideal}{理想}由关系 (1.13.1) 生成。因此可以应用第 V 讲 4.12 和 4.13 ii)，只须证明对每个 $n\geq0$，
\[
x^n\in\operatorname{Fil}^nK'(S)+\operatorname{Fil}^{n-1}K'(S)\,x+\cdots+\operatorname{Fil}^{n-r}K'(S)\,x^r.
\]
$n\leq r$ 时这显然成立；以下用归纳法。由关系 (1.13.1)，
\[
x^n=-\sum_{i=0}^r\gamma^{r+1-i}(\check E-r-1)\,x^{n-r-1+i},\qquad n\geq r+1.
\]
将归纳假设用于各 $x^{n-r-1+i}$，即得所求。由于 $1,\ldots,x^r$ 在 $K'(S)$ 上\Term{linearly-independent}{线性无关}，上述和是\Term{direct-sum}{直和}。
\end{proof}

\SourcePage{414}
\begin{SGAresult}{推论}{5.4}
$K'(X)$ 在 $K'(S)$ 上诱导的\Term{filtration}{滤过}与 $K'(S)$ 自身的\Term{filtration}{滤过}相同。
\end{SGAresult}

\begin{SGAresult}{推论}{5.5}
设 $S$ 为\Term{quasi-compact}{拟紧}\Term{scheme}{概形}，$\mathcal E_1,\ldots,\mathcal E_n$ 为\Term{finite-type}{有限型}\Term{locally-free-sheaf}{局部自由} $\mathcal O_S$-\Term{module}{模}，$\mathrm D(\mathcal E_i)$ 为 $\mathcal E_i$ 的\Term{flag-scheme}{旗概形}，$X=\mathrm D(\mathcal E_1)\times_S\cdots\times_S\mathrm D(\mathcal E_n)$。则\Term{canonical}{典范}\Term{homomorphism}{同态}
\[
\operatorname{Gr}K'(S)\longrightarrow\operatorname{Gr}K'(X)
\]
为\Term{injectivity}{单射}。
\end{SGAresult}
\begin{proof}
3.1 的论证可将问题归结为 $X$ 为 $S$ 上\Term{projective-bundle}{射影丛}时的同一性质。而这一情形直接来自 5.4。
\end{proof}

\begin{SGAresult}{推论}{5.6}
在 5.1 的记号与假设下，$\operatorname{Gr}K'(X)$ 是自由 $\operatorname{Gr}K'(S)$-\Term{module}{模}，以 $1,\bar x,\ldots,\bar x^r$ 为基。对任意 $k$，有
\[
\operatorname{Gr}^kK'(X)=\operatorname{Gr}^kK'(S)\oplus\operatorname{Gr}^{k-1}K'(S)\,\bar x\oplus\cdots\oplus\operatorname{Gr}^{k-r}K'(S)\,\bar x^r.
\]
最后，元素 $1,\bar x,\ldots,\bar x^{r+1}$ 满足关系
\[
\sum_{i=0}^{r+1}c^{r+1-i}(\check E)\bar x^i=0,
\]
其中 $c^k(\check E)$ 为 $\check E$ 的第 $k$ 个\Term{chern-class}{Chern 类}（第 V 讲 6.7）。
\end{SGAresult}
\begin{proof}
由 1.1、5.3 及 (1.13.1) 直接得到。
\end{proof}

\begin{SGAresult}{推论}{5.7}
在 4.5 的记号与假设下，$\operatorname{Gr}K'(S)$-\Term{algebra}{代数} $\operatorname{Gr}K'(X)$ 由 $\bar c_i^{(j)}$ 生成，其中 $\bar c_i^{(j)}$ 为 $\gamma^i(F_j-p_j)$ 在 $\operatorname{Gr}^iK'(X)$ 中的类。各 $\bar c_i^{(j)}$ 之间的关系\Term{ideal}{理想}由下式\Term{pushout}{推出}的关系生成：
\[
c(E-n)=\prod_{j=1}^k c(F_j-p_j),
\]
\SourcePage{415}
其中 $c$ 为\Term{total-chern-class}{总 Chern 类}（第 V 讲 6.7）。
\end{SGAresult}
\begin{proof}
重复 4.6 的证明，并使用 5.6 即可证明这一结果（比较 [2]，第 4 讲）。
\end{proof}

\begin{SGAresult}{推论}{5.8}
在 5.1 的记号与假设下，对任意 $k$，有
\[
f_*(\operatorname{Fil}^kK'(X))\subset\operatorname{Fil}^{k-r}K'(S).
\]
\end{SGAresult}
\begin{proof}
这是 5.3 和\Term{projection-formula}{投影公式}的直接推论。
\end{proof}

\begin{SGAresult}{命题}{5.9}
在 5.1 的记号与假设下，对每个 $i\geq1$，有
\[
f_*(\lambda^i(F))=0\qquad\text{（从而 }f_*(\sigma^i(F))=0\text{）}.
\]
特别地，置 $\lambda_{-1}(F)=\sum_{i=0}^r(-1)^i\lambda^i(F)$，则
\[
f_*(\lambda_{-1}(F))=1.
\]
\end{SGAresult}
\begin{proof}
定义一个映射
\[
\widetilde f_*\colon1+K'(X)[[t]]^+\longrightarrow1+K'(S)[[t]]^+,
\]
令
\[
\widetilde f_*\left(\sum_{i\geq0}a_it^i\right)=\sum_{i\geq0}f_*(a_i)t^i.
\]
它\Term{generic}{一般}不是群\Term{homomorphism}{同态}；不过，对 $\varphi\in1+K'(X)[[t]]^+$ 和 $\psi\in1+K'(S)[[t]]^+$，有如下关系：
\[
\widetilde f_*\bigl(\varphi\,\widetilde{g^*}(\psi)\bigr)=\psi\,\widetilde f_*(\varphi).
\]
事实上，可以逐个系数验证，而它直接来自\Term{projection-formula}{投影公式}。

于是，由关系 $E=F+\xi$ \Term{pushout}{推出}
\SourcePage{416}
\[
\widetilde f_*(\lambda_t(F))
=\widetilde f_*(\lambda_t(E)\lambda_t(-\xi))
=\lambda_t(E)\widetilde f_*(\lambda_t(-\xi)).
\]
而由 (4.2.3)，
\[
\widetilde f_*(\lambda_t(-\xi))=\sum_{i\geq0}(-1)^if_*(\xi^i)t^i=\sigma_{-t}(E).
\]
利用关系 (1.12.1)，最终得到
\[
\widetilde f_*(\lambda_t(F))=1,
\]
所求关系随之得证。
\end{proof}

\begin{SGAresult}{命题}{5.10}
在 5.1 的记号与假设下，设 $z\in K'(X)$。$z(1-\xi)=0$ 的充要条件是存在 $y\in K'(S)$，使
\[
\textup{i)}\quad z=y\lambda_{-1}(F);\qquad
\textup{ii)}\quad y\lambda_{-1}(E)=0.
\]
\end{SGAresult}
\begin{proof}
由 $E=F+\xi$ \Term{pushout}{推出}
\[
\lambda_{-1}(E)=\lambda_{-1}(F)\lambda_{-1}(\xi)=\lambda_{-1}(F)(1-\xi).
\]
因此条件是充分的。

设 $z(1-\xi)=0$，即 $zx'=0$。由注记 1.13，可以唯一地找到 $a_i\in K'(S)$，使
\[
z=\sum_{i=0}^r a_i x'^i.
\]
另一方面，各 $x'^i$ 满足一个\Term{linear}{线性}关系
\[
\sum_{i=0}^r c_i x'^i+x'^{r+1}=0,\qquad
c_i=(-1)^{r+i+1}\gamma^{r+1-i}(E-r-1).
\]
由于 $zx'=0$，得到
\[
a_rx'^{r+1}=-\sum_{i=1}^r a_{i-1}x'^i=-\sum_{i=0}^r a_rc_i x'^i.
\]
利用 $1,\ldots,x'^r$ 的\Term{linearly-independent}{线性无关}性，得
\[
a_rc_0=0;\qquad a_i=a_rc_{i+1}\quad(0\leq i<r).
\]
\SourcePage{417}
故
\[
z=a_r\left(\sum_{i=0}^{r-1}c_{i+1}x'^i+x'^r\right).
\]
而
\[
\begin{aligned}
\sum_{i=0}^{r-1}c_{i+1}x'^i+x'^r
&=(-1)^r\left(\sum_{i=0}^{r-1}\gamma^{r-i}(E-r-1)(-1)^ix'^i+(-1)^rx'^r\right)\\
&=(-1)^r\left(\sum_{i=0}^r\gamma^{r-i}(E-r-1)\gamma^i(-x')\right)\\
&=(-1)^r\gamma^r(E-r-1-x')=(-1)^r\gamma^r(F-x')\\
&=(-1)^r\lambda^r(F-1)\\
&=\lambda_{-1}(F).
\end{aligned}
\]
置 $a_r=y$，即得 $z=y\lambda_{-1}(F)$。此外，$yc_0=0$，而
\[
c_0=(-1)^{r+1}\gamma^{r+1}(E-r-1)=\lambda_{-1}(E),
\]
故 $y\lambda_{-1}(E)=0$。
\end{proof}

\begin{SGAresult}{命题}{5.11}
在 5.1 的记号与假设下，设 $z\in\operatorname{Gr}K'(X)$。$z\bar x'=0$ 的充要条件是存在 $y\in\operatorname{Gr}K'(S)$，使
\[
\textup{i)}\quad z=yc^r(F);\qquad
\textup{ii)}\quad yc^{r+1}(E)=0.
\]
\end{SGAresult}
\begin{proof}
在\Term{associated-graded-ring}{相伴分次环}中使用 5.6，证明与 5.10 完全相同。
\end{proof}

\SGAsection{6}{当 X 有丰沛层时其 $K'$ 的滤过的研究}
本节中，$X$ 为具有\Term{ample-sheaf}{丰沛层} $\mathcal L$ 的\Term{scheme}{概形}。回顾一下，此时由\Term{finite-type}{有限型}\Term{locally-free-sheaf}{局部自由层}或由\Term{perfect-complex}{完美复形}定义的两个 $K'$ 相等（第 IV 讲 2.9）。
\SourcePage{418}
以 $\varepsilon$ 表示\Term{augmentation}{增广}\Term{homomorphism}{同态} $K'(X)\longrightarrow H^0(X,\mathbf Z)$。

\begin{SGAresult}{命题}{6.1}
设 $X$ 为具有\Term{ample-sheaf}{丰沛层} $\mathcal L$ 的\Term{quasi-compact}{拟紧}\Term{scheme}{概形}。则 $\operatorname{Fil}^1K'(X)$ 是一个\Term{nil-ideal}{幂零元理想}（nilidéal）。
\end{SGAresult}
\begin{proof}
由于 $\operatorname{Fil}^1K'(X)$ 由形如 $E-\varepsilon(E)$ 的元素生成，其中 $E$ 为\Term{finite-type}{有限型}\Term{locally-free-sheaf}{局部自由} $\mathcal O_X$-\Term{module}{模} $\mathcal E$ 的类，$\varepsilon\colon K'(X)\longrightarrow H^0(X,\mathbf Z)$ 为 $K'(X)$ 的\Term{augmentation}{增广}，故只须证明这样的元素\Term{nilpotent}{幂零}。由 $X$ 的\Term{quasi-compact}{拟紧}性，以及 $K'$ 与\Term{finite}{有限}\Term{disjoint-union}{不交并}的相容性，可以归结为 $\mathcal E$ 为常\Term{rank}{秩}的情形。最后，转到 $\mathcal E$ 的\Term{flag-bundle}{旗丛}上，这在 $K'$ 上给出\Term{injectivity}{单射}（3.1），且保留存在\Term{ample-sheaf}{丰沛层}的假设，故可设 $\mathcal E$ \Term{invertible-sheaf}{可逆}。

设 $\mathcal E^*$ 为 $\mathcal E$ 的\Term{dual}{对偶}。对充分大的 $n$，$\mathcal E^*(n)=\mathcal E^*\otimes_{\mathcal O_X}\mathcal L^{\otimes n}$ 由其\Term{global-section}{整体截面}生成；因此可找到整数 $m$ 以及满\Term{homomorphism}{同态} $(\mathcal O_X)^m\longrightarrow\mathcal E^*(n)$，从而（与 $\mathcal E(-n)$ 作\Term{tensor-product}{张量积}）得到满\Term{homomorphism}{同态} $(\mathcal E(-n))^m\longrightarrow\mathcal O_X$。此\Term{homomorphism}{同态}的\Term{kernel}{核}为\Term{rank}{秩} $m-1$ 的\Term{locally-free-sheaf}{局部自由层}。以 $L$ 表示 $\mathcal L$ 的类，即得
\[
\lambda^m(mEL^{-n}-1)=0.
\]
而
\[
\lambda^m(mEL^{-n}-1)=(-1)^m\lambda_{-1}(mEL^{-n})
=(-1)^m\lambda_{-1}(EL^{-n})^m=(-1)^m(1-EL^{-n})^m.
\]
因此，对每个\Term{invertible-sheaf}{可逆层} $\mathcal E$，当 $n$ 充分大时，元素 $1-EL^{-n}$ \Term{nilpotent}{幂零}；或者说，它与 $L^n$ 的乘积 $L^n-E$ \Term{nilpotent}{幂零}。

于是可写成
\[
E-1=(L^n-1)-(L^n-E).
\]
由上面的结论，若 $n$ 选得充分大，则 $L^n-E$ 与 $L^n-1$ 均\Term{nilpotent}{幂零}，故 $E-1$ \Term{nilpotent}{幂零}。
\end{proof}

\SourcePage{419}
\begin{SGAresult}{推论}{6.2}
在 6.1 的假设下，设 $R\subset K'(X)$ 为由\Term{finite}{有限}个\Term{sheaf}{层}的类生成的\Term{lambda-ring}{$\lambda$-子环}。则 $R$ 的\Term{lambda-ring}{$\lambda$-环}\Term{filtration}{滤过}是\Term{discrete-filtration}{离散的}，且 $R\cap\operatorname{Fil}^1K'(X)$ 是\Term{nilpotent-ideal}{幂零理想}。此外，$R$ 是 $R\cap H^0(X,\mathbf Z)$ 上的\Term{finite-type}{有限型}\Term{module}{模}（第 V 讲 3.9.1）。
\end{SGAresult}
\begin{proof}
由于 $X$ \Term{quasi-compact}{拟紧}，可用\Term{finite}{有限}个两两不交的\Term{open-subset}{开集}\Term{covering}{覆盖} $X$，使得其类生成 $R$ 的那些\Term{sheaf}{层}在每个\Term{open-subset}{开集}上均为常\Term{rank}{秩}。通常的论证因而可将问题归结为这些\Term{sheaf}{层}在 $X$ 上具有常\Term{rank}{秩}的情形；这时 $R$ 的\Term{augmentation}{增广}取值于 $\mathbf Z$。

设 $\mathcal E_i$（$i=1,\ldots,k$）为其类生成 $R$ 的那些\Term{sheaf}{层}，$E_i$ 为它们的类，$x_i=E_i-\varepsilon(E_i)$。可将 $R$ 写成由 $x_i$ 在关系 $\gamma^j(x_i)=0$（$j>r_i=\varepsilon(E_i)$）下生成的\Term{lambda-ring}{$\lambda$-环}的\Term{quotient}{商}，也就是
\[
\mathbf Z[x_i,\gamma^j(x_i)]_{\substack{1\leq i\leq k\\2\leq j\leq r_i}}
\]
的\Term{quotient}{商}；因而 $R$ 为\Term{noetherian}{诺特}环。但由命题 6.1，各 $x_i,\gamma^j(x_i)$ 均\Term{nilpotent}{幂零}。故 $R$ 的\Term{augmentation-ideal}{增广理想}\Term{nilpotent-ideal}{幂零}；而此\Term{ideal}{理想}恰为 $R\cap\operatorname{Fil}^1K'(X)$。

$\operatorname{Fil}^nR$ 的每个元素都可写成关于 $x_i,\gamma^j(x_i)$ 的\Term{polynomial}{多项式}，且没有权 $<n$ 的\Term{set-component}{分量}。设 $m$ 为使得存在某个 $i$ 满足 $\gamma^{m-1}(x_i)\ne0$ 的最大整数；设 $p$ 为使所有 $x_i,\gamma^j(x_i)$ 的 $p$ 次幂均为零的整数。当 $n>kp(1+\cdots+m)$ 时，立即可知 $\operatorname{Fil}^nR=0$；因而 $R$ 的\Term{filtration}{滤过}是\Term{discrete-filtration}{离散的}。

由于 $R$ 的\Term{filtration}{滤过}\Term{discrete-filtration}{离散}，关于 $x_i,\gamma^j(x_i)$ 的非零\Term{monomial}{单项式}只有\Term{finite}{有限}个。这些\Term{monomial}{单项式}构成 $R$ 作为 $\mathbf Z$-\Term{module}{模}的一组\Term{generator}{生成元}，证明完毕。
\end{proof}

\SourcePage{420}
\begin{SGAresult}{推论}{6.3}
在 6.1 的假设下，设 $A$ 为\Term{commutative-ring}{交换环}。元素 $y\in K'(X)\otimes_{\mathbf Z}A=K'(X)_A$ \Term{invertible}{可逆}的充要条件是 $z=(\varepsilon\otimes\operatorname{Id}_A)(y)$ \Term{invertible}{可逆}。特别地，若 $\mathcal E$ 是 $X$ 上\Term{nowhere-zero}{处处非零}的\Term{finite-type}{有限型}\Term{locally-free-sheaf}{局部自由层}，则它在 $K'(X)_{\mathbf Q}$ 中的类\Term{invertible}{可逆}。
\end{SGAresult}
\begin{proof}
$K'(X)\otimes_{\mathbf Z}A$ 的\Term{augmentation-ideal}{增广理想}是 $\operatorname{Fil}^1K'(X)\otimes_{\mathbf Z}A$，故由 6.1 是\Term{nil-ideal}{幂零元理想}。因此 $K'(X)\otimes_{\mathbf Z}A$ 中一个元素\Term{invertible}{可逆}的充要条件是它的\Term{augmentation}{增广}\Term{invertible}{可逆}。
\end{proof}

\begin{SGAresult}{推论}{6.4}
设 $X$ 为具有\Term{ample-sheaf}{丰沛层}的\Term{quasi-compact}{拟紧}\Term{scheme}{概形}，$f\colon X'\longrightarrow X$ 为\Term{proper}{固有}且\Term{perfect-morphism}{完美}的\Term{morphism}{态射}。再假定以下两个条件之一成立：

\textup{i)} $f$ 为\Term{projective-morphism}{射影态射}；

\textup{ii)} $X$ 为\Term{noetherian}{诺特}\Term{scheme}{概形}，

从而“\Term{finiteness-theorem}{有限性定理}”成立（第 III 讲 2.2）。给定一个环 $A$，并考虑“\Term{direct-image}{正像}”\Term{homomorphism}{同态}（第 IV 讲 2.11）：
\[
f_*\colon K'(X')\otimes_{\mathbf Z}A\longrightarrow K'(X)\otimes_{\mathbf Z}A.
\]
最后假定存在元素 $x'\in K'(X')\otimes_{\mathbf Z}A$，使 $\varepsilon(f_*(x'))\in H^0(X,A)$ \Term{invertible}{可逆}。则\Term{homomorphism}{同态}
\[
f^*\colon K'(X)\otimes_{\mathbf Z}A\longrightarrow K'(X')\otimes_{\mathbf Z}A
\]
为\Term{injectivity}{单射}。
\end{SGAresult}
\begin{proof}
写出\Term{projection-formula}{投影公式}
\[
f_*(f^*(x)x')=xf_*(x').
\]
而 $f_*(x')$ 的\Term{augmentation}{增广}\Term{invertible}{可逆}，故本身也\Term{invertible}{可逆}（6.3）。因此复合\Term{homomorphism}{同态}
\SourcePage{421}
\[
x\longmapsto f^*(x)\longmapsto f^*(x)x'\longmapsto f_*(f^*(x)x')
\]
是\Term{injectivity}{单射}，从而 $f^*$ 也是\Term{injectivity}{单射}。
\end{proof}

\textbf{6.5。} 设 $X$ 为\Term{dimension}{维数} $d$、具有\Term{ample-sheaf}{丰沛层} $\mathcal L$ 的\Term{noetherian}{诺特}\Term{scheme}{概形}。给 $K'(X)$ 配备如下定义的\Term{filtration}{滤过}（参见 [1]）：$x$ 属于 $K'(X)_i$，当且仅当对 $X$ 的每个\Term{closed-subset}{闭集} $Y$，均存在\Term{perfect-complex}{完美复形} $\mathcal M^\bullet(Y)$，使
\[
\begin{gathered}
\textup{i)}\quad \operatorname{cl}(\mathcal M^\bullet(Y))=x;\\
\textup{ii)}\quad \operatorname{codim}\bigl(\operatorname{Supp}(H^*(\mathcal M^\bullet(Y)))\cap Y,Y\bigr)\geq i,\qquad
H^*(\mathcal M^\bullet(Y))=\bigoplus_{i\geq0}H^i(\mathcal M^\bullet(Y)).
\end{gathered}
\]
如此定义的各 $K'(X)_i$ 是 $K'(X)$ 的加法\Term{subgroup}{子群}。事实上，设 $x,x'\in K'(X)_i$，对 $X$ 中一个给定\Term{closed-subset}{闭集} $Y$，$\mathcal M^\bullet,\mathcal M'^\bullet$ 为相应的\Term{complex}{复形}；则 $x+x'$ 是\Term{complex}{复形} $\mathcal M^\bullet\oplus\mathcal M'^\bullet$ 的类，且
\[
\begin{aligned}
\operatorname{Supp}(H^*(\mathcal M^\bullet\oplus\mathcal M'^\bullet))
&=\operatorname{Supp}(H^*(\mathcal M^\bullet)\oplus H^*(\mathcal M'^\bullet))\\
&=\operatorname{Supp}(H^*(\mathcal M^\bullet))\cup\operatorname{Supp}(H^*(\mathcal M'^\bullet)).
\end{aligned}
\]
由于\Term{finite}{有限}个\Term{closed-subset}{闭集}之并的\Term{codimension}{余维数}是各\Term{closed-subset}{闭集}\Term{codimension}{余维数}的\Term{infimum}{下确界}，确有 $x+x'\in K'(X)_i$。

\begin{SGAresult}{命题}{6.6}
$K'(X)$ 的\Term{filtration}{滤过} 5.5 是\Term{ring-filtration}{环滤过}，且 $K'(X)_{d+1}=0$。
\end{SGAresult}
\begin{proof}
设 $x\in K'(X)_i$，$x'\in K'(X)_j$，$Y$ 为 $X$ 的\Term{closed-subset}{闭集}。存在\Term{complex}{复形} $\mathcal M^\bullet$，使 $x=\operatorname{cl}(\mathcal M^\bullet)$，且
\[
\operatorname{codim}\bigl(\operatorname{Supp}(H^*(\mathcal M^\bullet))\cap Y,Y\bigr)\geq i;
\]
还存在\Term{complex}{复形} $\mathcal M'^\bullet$，使 $x'=\operatorname{cl}(\mathcal M'^\bullet)$，且
\[
\operatorname{codim}\bigl(\operatorname{Supp}(H^*(\mathcal M^\bullet))\cap\operatorname{Supp}(H^*(\mathcal M'^\bullet))\cap Y,
\operatorname{Supp}(H^*(\mathcal M^\bullet))\cap Y\bigr)\geq j.
\]
于是 $xx'=\operatorname{cl}(\mathcal M^\bullet\mathbin{\otimes^{\mathbf L}}\mathcal M'^\bullet)$，且
\[
\operatorname{codim}\bigl(\operatorname{Supp}(H^*(\mathcal M^\bullet))\cap\operatorname{Supp}(H^*(\mathcal M'^\bullet))\cap Y,Y\bigr)\geq i+j
\]
（EGA $0_{\mathrm{IV}}$ 14.2.2.3）。
\SourcePage{422}
因此，为完成证明，只须注意
\[
\operatorname{Supp}\bigl(H^*(\mathcal M^\bullet\mathbin{\otimes^{\mathbf L}}\mathcal M'^\bullet)\bigr)
\subset\operatorname{Supp}(H^*(\mathcal M^\bullet))\cap\operatorname{Supp}(H^*(\mathcal M'^\bullet)).
\]
为证明 $K'(X)_{d+1}$ 为零，取 $x\in K'(X)_{d+1}$，以及 $Y=X$；则存在 $\mathcal M^\bullet$，使 $x=\operatorname{cl}(\mathcal M^\bullet)$，且 $\operatorname{codim}_X\operatorname{Supp}(H^*(\mathcal M^\bullet))\geq d+1$。由于 $X$ 的\Term{dimension}{维数}为 $d$，这蕴涵 $\operatorname{Supp}(H^*(\mathcal M^\bullet))=\varnothing$，从而 $\mathcal M^\bullet$ \Term{acyclic-complex}{无环}，并有 $x=0$。
\end{proof}

\begin{SGAresult}{\Term{lemma}{引理}}{6.7}
设 $X=\operatorname{Spec}(A)$ 为\Term{affine-scheme}{仿射}\Term{noetherian}{诺特}\Term{scheme}{概形}，$x_1,\ldots,x_k$ 为 $X$ 的\Term{finite}{有限}个点，$\mathcal E$ 为\Term{invertible-sheaf}{可逆} $\mathcal O_X$-\Term{module}{模}。则存在一个\Term{section}{截面} $s\in\Gamma(X,\mathcal E)$，使对每个 $i$，$s_{x_i}$ 是 $\mathcal E_{x_i}$ 的一个基。
\end{SGAresult}
\begin{proof}
置 $M=\Gamma(X,\mathcal E)$。对每个 $i$，取 $X$ 的一个\Term{closed-point}{闭点} $x_i'$，使 $x_i'\in\overline{\{x_i\}}$。各 $x_i'$ 对应于 $A$ 的极大\Term{ideal}{理想} $\mathfrak m_i$；改变指标集，使各 $\mathfrak m_i$ 在族 $(\mathfrak m_i)$ 中恰出现一次，即可假设它们两两不同。由于 $\mathcal E$ \Term{invertible-sheaf}{可逆}，对每个 $i$ 有\Term{isomorphism}{同构} $M_{\mathfrak m_i}\simeq A_{\mathfrak m_i}$，从而有\Term{isomorphism}{同构}
\[
M/\mathfrak m_iM\simeq A/\mathfrak m_i,\qquad
\prod_iM/\mathfrak m_iM\simeq\prod_iA/\mathfrak m_i,
\]
亦即
\[
A/\prod_i\mathfrak m_i\simeq M/\bigl(\prod_i\mathfrak m_i\bigr)M.
\]
\Term{lifting}{提升}最后这个\Term{isomorphism}{同构}，得到一个\Term{homomorphism}{同态} $A\longrightarrow M$，使对每个 $i$，相应的\Term{homomorphism}{同态} $A_{\mathfrak m_i}\longrightarrow M_{\mathfrak m_i}$ 模 $\mathfrak m_i$ 后均为\Term{isomorphism}{同构}，故由 \Term{nakayama-lemma}{Nakayama 引理}本身即为\Term{isomorphism}{同构}。因此 $1$ 在这个\Term{homomorphism}{同态}下的像给出 $\mathcal E$ 的一个\Term{section}{截面} $s$，它是每个 $\mathcal E_{x_i'}$ 的基。由于 $s$ 为 $\mathcal E$ 的基的点构成 $X$ 的一个\Term{open-subset}{开集}，$s$ 也在 $x_i$ 处构成 $\mathcal E_{x_i}$ 的基，$i=1,\ldots,k$。
\end{proof}

\begin{SGAresult}{\Term{lemma}{引理}}{6.8}
设 $X$ 为具有\Term{ample-sheaf}{丰沛层} $\mathcal L$ 的\Term{noetherian}{诺特}\Term{scheme}{概形}，$\mathcal E$ 为\Term{invertible-sheaf}{可逆} $\mathcal O_X$-\Term{module}{模}，$x_1,\ldots,x_k$ 为 $X$ 的点。则存在一个整数 $n$，使对每个 $k\geq1$，均存在\Term{section}{截面} $s\in\Gamma(X,\mathcal E(kn))$，它对每个 $i$ 都是 $\mathcal E(kn)_{x_i}$ 的基。
\end{SGAresult}
\SourcePage{423}
\begin{proof}
可找到整数 $n_0$ 和一个\Term{section}{截面} $f\in\Gamma(X,\mathcal L^{\otimes n_0})$，使每个 $x_i$ 均属于 $X_f$，且 $X_f$ 为\Term{affine-scheme}{仿射}\Term{open-subset}{开集}（EGA II 4.5.4）。在 $X_f$ 上可找到满足\Term{lemma}{引理} 6.7 条件的\Term{section}{截面} $t\in\Gamma(X_f,\mathcal E)$；此外存在整数 $n_0'$，使\Term{section}{截面}
\[
f^{\otimes n_0'}\otimes t\in\Gamma(X_f,\mathcal E(n_0n_0'))
\]
延拓到整个 $X$。置 $n=n_0n_0'$；则 $f^{\otimes n_0'}\in\Gamma(X,\mathcal L^n)$；将上面所得 $\mathcal E(n)$ 的\Term{section}{截面}与 $f^{\otimes kn_0'}$ 作\Term{tensor-product}{张量积}，便对每个 $k\geq0$ 得到 $\mathcal E((k+1)n)$ 的一个\Term{section}{截面}。由于对一切 $p$，$X_{f^{\otimes p}}=X_f$，且 $t$ 在每个 $x_i$ 处为 $\mathcal E_{x_i}$ 的基，如此构造的 $\mathcal E(kn)$ 的\Term{section}{截面}对每个 $k\geq1$、每个 $i$，均构成 $\mathcal E(kn)_{x_i}$ 的基。
\end{proof}

\begin{SGAresult}{\Term{theorem}{定理}}{6.9}
设 $X$ 为\Term{dimension}{维数} $d$、具有\Term{ample-sheaf}{丰沛层} $\mathcal L$ 的\Term{noetherian}{诺特}\Term{scheme}{概形}。则 $\operatorname{Fil}^{d+1}K'(X)=0$，从而更有
\[
\bigl(\operatorname{Fil}^1K'(X)\bigr)^{d+1}=0.
\]
\end{SGAresult}
（在这里所考虑的\Term{generic}{一般}情形下，最后这个关系由 J.-P. Serre 得到。）
\begin{proof}
先证明：若 $\mathcal E,\mathcal E'$ 是 $X$ 上两个\Term{invertible-sheaf}{可逆层}，其类为 $E,E'$，则 $E-E'\in K'(X)_1$。

设 $Y$ 为 $X$ 的\Term{closed-subset}{闭集}，$Y_1,\ldots,Y_k$ 为其\Term{irreducible-component}{不可约分支}，并对每个 $i$ 取一点 $x_i\in Y_i-\bigcup_{j\ne i}Y_j$。由\Term{lemma}{引理} 6.8，存在整数 $n,n'$，使无论整数 $k,k'\geq1$ 如何选取，均可找到从 $\mathcal O_X$ 到 $\mathcal E(kn),\mathcal E'(k'n')$ 的\Term{homomorphism}{同态}，它们在各点 $x_i$ 处均为\Term{isomorphism}{同构}；由此得到\Term{morphism}{态射}
\[
\mathcal L^{\otimes(-kn)}\xrightarrow{s_k}\mathcal E,\qquad
\mathcal L^{\otimes(-k'n')}\xrightarrow{s'_{k'}}\mathcal E',
\]
\SourcePage{424}
它们在各点 $x_i$ 处均为\Term{isomorphism}{同构}。现考虑\Term{complex}{复形}
\[
\mathcal N^\bullet=\cdots\longrightarrow0\longrightarrow
\mathcal L^{\otimes(-nn')}\xrightarrow{(s'_n,0)}
\mathcal E'\oplus\mathcal L^{\otimes(-nn')}
\xrightarrow{0+s_{n'}}\mathcal E\longrightarrow0\longrightarrow\cdots.
\]
$\mathcal N^\bullet$ 的类为 $E-E'$；此外，$\mathcal N^\bullet$ 在各点 $x_i$ 处\Term{acyclic-complex}{无环}，故 $\operatorname{Supp}(H^*(\mathcal N^\bullet))\cap Y$ 不包含 $Y$ 的任何\Term{irreducible-component}{不可约分支}，从而在 $Y$ 中的\Term{codimension}{余维数}至少为 $1$；这表明 $E-E'$ 属于 $K'(X)_1$。

现设 $x\in\operatorname{Fil}^{d+1}K'(X)$，我们将证明 $x=0$。可假定 $x$ 写成
\[
x=\gamma^{i_1}(E_1-F_1)\cdots\gamma^{i_k}(E_k-F_k),\qquad i_1+\cdots+i_k\geq d+1,
\]
其中各 $E_i,F_i$ 为\Term{finite-type}{有限型}\Term{locally-free-sheaf}{局部自由层}的类，且对每个 $i$，$\mathcal E_i$ 与 $\mathcal F_i$ 的\Term{rank}{秩}相同。此外，将 $X$ \Term{resolution}{分解}为\Term{finite}{有限}个\Term{open-subset}{开集}的\Term{disjoint-union}{不交并}，并在每个\Term{open-subset}{开集}上分别论证，即可假设各 $\mathcal E_i,\mathcal F_i$ 均为常\Term{rank}{秩}。设 $X'$ 为各 $\mathcal E_i,\mathcal F_i$ 的\Term{flag-scheme}{旗概形}在 $X$ 上的\Term{fiber-product}{纤维积}。在 $K'(X')$ 中，$E_i,F_i$ 的像均为\Term{invertible-sheaf}{可逆层}之类的和。设 $\xi,\eta$ 为两个\Term{invertible-sheaf}{可逆} $\mathcal O_{X'}$-\Term{module}{模}的类；则
\[
\gamma_t(\xi-\eta)=\gamma_t(\xi-1)/\gamma_t(\eta-1)
=\frac{1+(\xi-1)t}{1+(\eta-1)t}.
\]
故 $\gamma^i(\xi-\eta)=(-1)^{i-1}(\xi-\eta)(\eta-1)^{i-1}$。而 $X'$ 是具有\Term{ample-sheaf}{丰沛层}的\Term{noetherian}{诺特}\Term{scheme}{概形}，故由上面的结论及 6.6，$\gamma^i(\xi-\eta)\in K'(X')_i$。再由 $\gamma_t$ 的\Term{additive}{加性}及 6.6，得到
\[
f^*\gamma^{i_j}(E_j-F_j)\in K'(X')_{i_j},\qquad
f^*(x)\in K'(X')_{d+1},
\]
其中 $f$ 为\Term{canonical}{典范}\Term{morphism}{态射} $X'\longrightarrow X$。

由此将\Term{pushout}{推出} $f^*(x)=0$。设 $r$ 为 $X'$ 在 $X$ 上的\Term{relative-dimension}{相对维数}。$K'(X')$ 中存在一个元素 $x'$，属于 $K'(X')$ 在 $K'(X)$ 上的一组基，且满足
\SourcePage{425}
\[
x'=\prod_{i=0}^r(\xi_i-1),
\]
其中各 $\xi_i$ 为\Term{invertible-sheaf}{可逆层}的类。为证明这一点，按 3.1 的论证可归结为 $X$ 上一个\Term{projective-bundle}{射影丛}的情形；而后一情形直接来自 1.13。由前面的结论，这个 $x'$ 属于 $K'(X')_r$，从而 $x'f^*(x)\in K'(X')_{d+r+1}$。而 $\dim X'=d+r$；故由 6.6，$x'f^*(x)=0$。由于 $x'$ 是 $K'(X')$ 在 $K'(X)$ 上的一组基的一个元素，这蕴涵 $f^*(x)=0$。最后，由 $f^*$ 是\Term{injectivity}{单射}（3.1），得到 $x=0$。
\end{proof}

\begin{SGAresult}{问题}{6.10}
在 6.9 的假设下，是否对每个 $k$ 都有
\[
\operatorname{Fil}^kK'(X)\subset K'(X)_k\ ?
\]
\end{SGAresult}

\bigskip
\noindent\textbf{参考文献}
\begin{enumerate}
\item[{[1]}] H. Bass，\textit{K 理论与\Term{stable}{稳定}\Term{algebra}{代数}}（K-Theory and stable algebra），IHÉS 数学出版物，第 22 号。
\item[{[2]}] C. Chevalley，\textit{1958 年讨论班：Chow 环及其应用}（Séminaire 1958 : Anneaux de Chow et applications）。
\item[{[3]}] M. Hakim，\textit{\Term{locally-ringed-topos}{局部环化拓扑斯}与\Term{relative-scheme}{相对概形}}（Topos localement annelés et schémas relatifs），学位论文（准备中）。
\item[{[4]}] R. Hartshorne，\textit{\Term{residue}{留数}与\Term{duality}{对偶性}}（Residues and duality），数学讲义丛书（Lecture Notes in Mathematics），第 20 卷，Springer-Verlag。
\item[{[5]}] J.-P. Serre，\textit{\Term{reductive-group-scheme}{约化群概形}的 Grothendieck 群}（Le groupe de Grothendieck des schémas en groupes réductifs），IHÉS 数学出版物，第 34 号。
\end{enumerate}
